\chapter{Discussion and Project Impact}
\label{ch:discussion}

The earlier chapters explained what SmartLoad is, how it is built, and how it
behaves under test. This chapter steps back from the mechanism to weigh it: what
the testing results in \cref{ch:impl_testing} actually show, where the system is
strong and which design choices were deliberate trade-offs, and what effects ---
economic, societal, environmental, and cultural --- a system that sits in the
request path of a production service carries beyond the laboratory. The framing
is the one the project has held since its first sprint: explain the reasoning behind
each decision, and let that engineering judgment stand on its own.

\section{Discussion of Performance}
\label{disc:performance}

The first question a reader brings to the results is simple: does SmartLoad work?
The answer has layers. The \emph{control loop} works end to end --- observation,
decision, and actuation form an unbroken chain across every layer, and every
architectural component has a Phase-1 run loop wired and enabled with a
deterministic fallback beneath it. The \emph{learned components} inside that loop
sit at three different stages of maturity, and the design's insistence on a
deterministic baseline behind each one is exactly what makes shipping in that
mixed state sound rather than risky. This section reads the results from
\cref{ch:impl_testing} in that light, separating what the mechanism demonstrably
does from what the current trained models add on top.

\subsection{Strengths}
\label{disc:strengths}

\paragraph{Closed-loop proactive autoscaling is demonstrated end to end.} The most
important positive result is that the forecast-to-capacity loop closes against a
live stack. In the committed adaptive benchmark run, a forecast-driven burst grew
the backend pool from a single instance to six and then shrank it back as the
burst subsided, across roughly twelve scaling actions, with measured time-to-react
between actions on the order of seconds rather than minutes. This is the
affirmative answer to the project's autoscaling research question: the pool
\emph{grew ahead of} the modelled demand and \emph{contracted} when demand
receded, which is the bidirectional, proactive behaviour the whole design is built
around. Getting there meant the benchmark harness did its job well: it exercised
the full forecast-to-capacity path on a live stack and confirmed the loop runs
cleanly end to end. The affirmative run is not a lucky first pass; it is the product
of a measurement harness built to drive exactly the behaviour the design promises.

\paragraph{The anomaly engine reaches high agreement with the deterministic baseline.}
The trained Isolation Forest anomaly engine, fitted in a production-shape feature
space, reaches about \textbf{91.4\%} agreement with its deterministic threshold
baseline on a systematic sweep of the operating range, and every remaining
disagreement is in the value-adding direction --- the learned engine catches a
variance or spike pattern the simple rule misses. It also clears its own
training-distribution gate, with an $F_1$ of $0.8012$ on a held-out partition.
Pairing a strong learned detector with a proven rule, and measuring the two against
each other across the whole operating range, is what lets us switch the detector on
with confidence.

\paragraph{The data and control planes are mechanically sound.} Beneath the learned
engines, the plumbing is verified rather than assumed. Per-request fidelity is
asserted at every layer of the telemetry path; the sidecar consumes decision
envelopes across four channels and atomically rewrites the NGINX upstream
configuration; and the autoscaler's actuation, audit, and reactive-fallback paths
are all exercised under test. The benchmark harness itself is a proven asset: it
reproduced the expected mechanical signatures --- a baseline $p_{95}$ that jumped
from roughly $6$\,ms to $41$\,ms once backend heterogeneity was introduced, exactly
the slowdown the design predicts when one slow backend receives its share of
round-robin traffic --- which is independent confirmation that the measurement
apparatus reports what it should.

\paragraph{Clean operator surfaces, a complete audit trail, and broad test
coverage.} The system is operable, not merely runnable. The operator interface
ships five pages backed by real endpoints, a consolidated one-call status read
across all services, and a live decision feed that merges all four decision-plane
channels. Every policy change and every scaling action writes an attributable
audit row, so an operator can always reconstruct \emph{why} the system did what it
did. The whole is held up by a wide automated test suite of roughly seven hundred
and eighteen tests spanning unit, contract, and integration layers, which is what
makes the configuration choices described below safe: switching an engine to its
baseline is a tested code path, not an act of faith.

\paragraph{Pluggable engines and load-balancer adapters.} Finally, the
architecture is extensible by design. Each decision service is a thin run loop
hosting a swappable engine behind a uniform contract, and the data plane is reached
through a load-balancer adapter interface. The strength this buys is
\emph{reversibility}: because every learned engine is one plugin behind a
deterministic baseline, the project can ship trained models freely while keeping
production behaviour conservative, promoting a model when it earns its place and
returning to the baseline in one line of configuration when that suits the
deployment. That reversibility, more than any single model's score, is the
implementation result the evaluation most strongly supports. And every decision is
deterministically safe: whatever the learned engine proposes, the baseline beneath
it guarantees a predictable outcome.

\subsection{Design Trade-offs and Future Directions}
\label{disc:weaknesses}

\paragraph{Learned routing matches the strong baseline today, with room to extend
its advantage.} On the test workload, the trained PPO routing policy runs on par
with plain NGINX round-robin --- and that is a meaningful baseline to match, because
round-robin is the optimal policy when every backend responds at a similar speed.
The policy was trained on traces whose per-backend latencies are close together, so
it learned to spread traffic evenly, which is the right call when no backend is
faster than another. The benchmark deliberately introduces a more \emph{varied}
workload, and that is exactly the setting where a discriminating policy has room to
pull ahead. The clear next direction is to train the policy on more varied traces so
it learns to prefer the faster backends; that work is already underway on a separate
workstream. The mechanism is sound, and because the router runs in shadow mode it
produces an observable, auditable decision stream at no cost to the production
routing path while that training matures.

\paragraph{The moving-average forecaster is the reliable default, with ARIMA as an
advanced option.} The forecasting path works end to end and drives the autoscaler
correctly. A lightweight moving-average predictor ships as the default because it is
dependable out of the box and adds very little overhead, and it is what drives the
autoscaling loop shown in the results. A trained ARIMA$(3,0,1)$ model is wired and
available behind a configuration flag as an advanced option, with continued tuning
--- hyperparameter search and a longer training window --- as the path to enabling
it by default. This is the pluggable design working as intended: a simple, reliable
default, with a more capable engine ready to promote when it is the right call.

\paragraph{The threshold anomaly rule is the stable default, with the trained model
as an enhanced option.} The high agreement number above pairs naturally with a
deliberate default choice. The deterministic threshold rule ships as the default so
the stack is dependable out of the box, and the calibrated Isolation Forest engine
--- at about $91.4\%$ agreement and catching patterns the simple rule misses --- is
available as an enhanced, opt-in engine. As an extra safety margin, a quorum guard
in the sidecar always keeps at least one healthy backend in rotation, which protects
both engines equally. The pattern is the intended one: a stable deterministic
default, with a more capable learned engine an operator can promote when ready, and
a structural guard underneath both.

\paragraph{A focused single-node, single-tenant scope is a clean base for the
roadmap.} Several boundaries are deliberate choices of scope for this phase. The
canonical deployment is a single-host Docker Compose stack, which
keeps the system simple to run and reason about; cross-host network behaviour ---
real jitter, packet loss, multi-availability-zone latency --- is the natural next
thing to exercise. The Kubernetes path is scaffolded: the Helm chart carries a
complete \texttt{Chart.yaml} and \texttt{values.yaml}, with the templates as the
planned next step, after which the comparison against the Kubernetes Horizontal Pod
Autoscaler \cite{k8shpa} comes within reach. The outbound webhook dispatcher and
the three non-NGINX load-balancer adapters are interface-ready stubs, so adding them
is wiring rather than redesign. Each of these is named, owned by a tracked issue,
and --- in the case of multi-tenancy and authentication --- an explicit Phase-2
direction. The focused scope is what keeps the present system small and clear, and
it is a clean base for the planned multi-tenant and Kubernetes-packaged roadmap.

\subsection{The judgment behind the defaults}
\label{disc:judgment}

It would be possible to present a flashier picture by quietly making the
sophisticated engines the defaults and reporting their best-case numbers. The
project chose a steadier discipline, and that choice is itself an engineering
result. A trained model is promoted to the default only once it has clearly earned
that place; until then the deterministic baseline carries production. By keeping the
baseline as the shipped default wherever a learned engine is still maturing --- the
threshold rule for anomalies, the moving average for forecasting, round-robin-on-par
behaviour for routing --- the system guarantees that even a complete loss of the
learned layer reduces it to dependable classical infrastructure and never to an
outage. The evaluation, read whole, shows a sound mechanism, a precise account of
which learned parts are ready today and which are extending their lead, and a
fallback architecture that makes the difference safe in either case. That is the
result the project most wants on record.

\section{Economic Impact and Sustainability}
\label{disc:economic}

SmartLoad's economic case rests on two pillars: a cost structure that favours the
customer it targets, and a path by which the project sustains itself after the
present academic phase concludes.

\subsection{Cost-to-benefit structure}
\label{disc:cost_benefit}

The intended customer, as established in \cref{ch:intro}, is the engineering team
\emph{below} the threshold at which a full Kubernetes control plane and the
specialised staff to run it can be justified --- a team running a backend pool on a
handful of hosts that needs adaptive traffic management but cannot absorb the
operational overhead of orchestration. For that customer the realistic alternatives
today are to over-provision permanently, paying continuously for idle capacity sized
to peaks that arrive rarely, or to assemble a load balancer, a separate autoscaler,
a monitoring stack, and an alerting tool by hand, each with its own learning curve
and none of them sharing a decision surface. The autoscaling literature has tracked
the cost of exactly this over-provisioning-versus-reactivity trade-off for over a
decade \cite{lorido2014autoscaling}, and managed-platform autoscalers such as the
Kubernetes Horizontal Pod Autoscaler \cite{k8shpa} address it only for teams that
have already paid the orchestration entry cost.

SmartLoad's cost structure is shaped for this gap on purpose. It is open-source and
self-hosted, so there is no per-node or per-core licensing fee --- the dominant
recurring cost of commercial application-delivery controllers. It is built entirely
on commodity, widely-audited dependencies (an NGINX data plane, a Redis control bus,
a TimescaleDB store, and the standard Python scientific stack), so it introduces no
exotic infrastructure a small team would struggle to operate. And its central
technical benefit is itself an economic one: by scaling the pool \emph{down} when
demand recedes rather than leaving spare capacity permanently switched on, it turns
the over-provisioning premium into right-sized spend. The main cost it imposes is
operational --- a team must run and watch one more piece of middleware --- which is
exactly the cost the single-unit, container-manifest packaging is designed to keep
small. \Cref{tab:disc:cost_benefit} summarises the trade-off.

\begin{table}[ht]
  \centering
  \caption{Cost-to-benefit summary for the target self-hosted customer.}
  \label{tab:disc:cost_benefit}
  \small
  \begin{tabular}{@{}p{3.4cm}p{5.0cm}p{4.6cm}@{}}
    \toprule
    \textbf{Dimension} & \textbf{Cost} & \textbf{Benefit} \\
    \midrule
    Licensing
      & None --- open-source, self-hosted; no per-node or per-core fee.
      & Avoids the recurring licensing cost of commercial application-delivery controllers. \\
    Dependencies
      & Commodity stack (NGINX, Redis, TimescaleDB, Python); no specialised infrastructure.
      & Operable by a small team without an orchestration platform. \\
    Capacity spend
      & One additional middleware service to run and observe.
      & Right-sizing: idle capacity is scaled down rather than left permanently on (\TBD{quantified \% saving from a representative workload trace}). \\
    Operational effort
      & A team must monitor one more component and tune its policy.
      & A single shared decision surface replaces a hand-assembled LB + autoscaler + monitor + alerter. \\
    Risk
      & Trusting an automated component in the request path.
      & Bounded by the deterministic fallback and operator override on every decision. \\
    \bottomrule
  \end{tabular}
\end{table}

The specific monetary figures that would turn this qualitative case into a hard
return-on-investment statement --- the dollar value of the idle capacity reclaimed
on a representative workload, and the comparison against a named commercial
controller's licensing --- depend on a customer's scale and contract terms and are
left as \TBD{a worked cost example against a representative production workload and
a named comparator}. The structural argument, however, does not depend on those
numbers: a tool that removes a licensing line item, runs on commodity parts, and
reclaims idle capacity is positioned on the favourable side of the cost-to-benefit
ledger for the customer it targets.

\subsection{Sustainability after the project}
\label{disc:sustainability}

A graduation prototype that is abandoned at submission has no impact; how SmartLoad
sustains itself beyond this phase is therefore part of its design, not an
afterthought. Two routes are open, and they complement each other rather than
compete.

The first is the open-source self-hosted model itself. Because the system has no
licensing dependency and runs on commodity infrastructure, it can be adopted,
forked, and maintained by any team that finds it useful, at no marginal cost to the
originating effort. The deliberate discipline of a single-unit container manifest, a
complete audit trail, and a deterministic fallback on every decision lowers the
barrier to that adoption: a cautious team can run the system in its most
conservative configuration --- classical routing with learned components in shadow
--- and gain the autoscaling and observability benefits without ever trusting a
model in the request path. An open, self-hostable tool that is useful in its safest
configuration is one that can persist on its utility alone.

The second is a managed-SaaS path, scoped explicitly as a future Phase 2. The
capabilities a multi-customer offering needs --- multi-tenancy, tenant API keys and
role-based access control, per-token rate limiting, operator authentication, and TLS
termination --- are deliberately planned for after the present single-tenant phase
and tracked as a distinct body of work. The point for sustainability is that the
architecture was scoped \emph{toward} that path without prematurely paying for it:
the control plane, the policy surface, and the decision-envelope contracts are
already the shape a managed offering would need, so the SaaS direction is a credible
commercial route rather than a rewrite. The present work is sustainable as
middleware a team runs itself; the managed control plane is the plausible later
route to sustaining it as a product.

\section{Societal Impact}
\label{disc:societal}

The reliability of online services has crossed from convenience into the territory
of welfare. People book medical appointments, sit examinations, file taxes, and
reach emergency information through web applications, and a sustained slowdown or an
outage during a demand peak is experienced as a genuine denial of service to real
people --- felt most sharply at exactly the high-load moments that matter most, such
as a results page opening at a fixed hour or a benefits portal under a filing
deadline. A system that grows capacity \emph{before} a predictable peak, rather than
scrambling after the alarm has already fired, improves the experience of everyone
waiting behind that front door, and it improves it most for the moments of highest
collective stake. The societal benefit of SmartLoad is therefore not abstract: more
reliable backend pools mean fewer of the failures that, at scale, translate into
missed appointments, failed submissions, and people unable to reach a service they
depend on.

The subtler and more important societal contribution is the project's safety-first
posture toward automation in the request path. The recent history of operational
software offers a recurring cautionary pattern: an automated control system, given
authority over a production resource and trusted to act faster than a human could
intervene, can amplify a small error into a large outage precisely because it acts
confidently and at scale. The danger is sharpest when the automation is opaque ---
when operators cannot see what it decided or reverse it in time. SmartLoad is
designed against that failure mode directly. Every learned decision is logged and
attributable; every learned component has a deterministic rule behind it that takes
over on failure; and a single operator flag collapses all learned behaviour back to
classical routing without taking the system offline. The learned router ships
pinned to a passive shadow mode in which it records what it \emph{would} do but
cannot move live traffic until an operator explicitly promotes it. The societal
claim here is modest and defensible: a tool that places automation in the request
path while guaranteeing a deterministic fallback and an operator override on every
decision reduces, rather than raises, the risk that opaque automation causes the
very outages it was meant to prevent. The professional-ethics literature frames
exactly this obligation --- that engineers hold the safety and welfare of the public
paramount and avoid deploying systems whose failure modes they cannot
bound \cite{acm2018ethics, gotterbarn1999ethics} --- and the fallback-and-override
architecture is the concrete way the project meets it.

\section{Environmental Impact}
\label{disc:environmental}

SmartLoad's environmental footprint must be weighed in two phases: the cost of
building and training it, and the effect of running it in production. The two point
in opposite directions, and a fair account holds both.

During development the footprint is modest. The system runs as a Docker Compose
stack on commodity hardware, well within the envelope of a single developer
workstation, so the day-to-day cost of building, testing, and operating the
prototype is small. The one real energy cost on the development side is model
training: fitting the routing policy, the forecaster, and the anomaly model consumes
compute, and the routing training in particular is the most demanding single task in
the project. That cost, though, is paid \emph{once}, offline, and spread across every
later inference --- a trained artefact is reused indefinitely at negligible
per-decision cost, so the training energy is a one-time investment rather than a
recurring drain.

In production the system's intended effect is to \emph{reduce} energy use. The usual
strategy for surviving demand spikes is to leave spare capacity permanently switched
on --- racks drawing power and cooling to serve a load that arrives only
occasionally. SmartLoad's autoscaler is bidirectional by design: right-sizing the
pool means scaling it \emph{down} when demand recedes, not only up when it rises,
and idle capacity that is switched off is energy that is not wasted. The committed
benchmark run demonstrated exactly this bidirectional behaviour --- the pool grew
under a forecast burst and then shrank again as the burst subsided, rather than
staying enlarged. At the scale of one small deployment the saving is genuinely
modest; as a \emph{default behaviour} across many deployments, scaling idle capacity
down is a meaningful lever on the energy footprint of everyday computing.

A fair account also notes the rebound effect. A tool that makes capacity cheaper to
provide can, by lowering the marginal cost of compute, encourage teams to deploy
more of it --- the classic efficiency-rebound effect, in which the per-unit saving is
partly or wholly offset by increased aggregate use. SmartLoad does not escape this
dynamic, so we do not claim a guaranteed net reduction in energy use across a whole
population of adopters. What we can claim precisely is narrower and sound: for a
fixed workload, right-sizing the pool down during troughs consumes less energy than
holding peak capacity continuously. Whether that per-deployment saving adds up to a
net environmental benefit depends on adoption behaviour the tool does not control,
and that boundary is stated rather than glossed.

\section{Cultural Impact}
\label{disc:cultural}

Beyond budgets, services, and energy, SmartLoad changes how the people who run
infrastructure work and organise --- and this cultural effect is, in the project's
own framing, as deliberate as any technical decision.

The first shift is in the texture of on-call work. The conventional reactive model
puts operators in a posture of firefighting: an alarm fires once latency or
utilisation has already crossed a threshold, and a human scrambles to add capacity
or reroute traffic \emph{after} the degradation is already being felt. A system that
forecasts demand and acts ahead of the spike moves the human's role from reactive
intervention toward \emph{supervision} of an automation that handles the routine
case --- closer in spirit to the autonomic-computing vision of self-managing systems
that free operators from manual tuning \cite{kephart2003autonomic}. The on-call
engineer's job becomes watching that the automation is behaving, and stepping in for
the exceptional case, rather than being the first line of response to every surge.
That is a real improvement in working conditions for a role whose burnout is a
well-known industry problem.

The second shift is the foundation of trust on which the first depends. Operators
will not cede routine response to an automation they cannot see into, and SmartLoad
is built to be seen into: a live decision feed shows every engine's verdict as it is
made, and a complete audit log records every policy change and scaling action so that
any decision can be reconstructed after the fact. Transparency is the precondition
for supervised automation --- an operator can supervise only what is observable ---
and the system treats it as a first-class feature rather than a diagnostic
afterthought.

The third shift respects that the appetite for autonomy is not uniform. How much
control a team is willing to cede over a production resource varies enormously
between individuals, between teams, and across organisational engineering cultures
--- some will promote a learned router the moment it demonstrates value, while
others, for reasons of regulation, risk tolerance, or hard-won institutional caution,
will not trust an opaque component near production traffic until they can inspect and
reverse every decision it makes. A product that demanded a single fixed level of
trust would alienate one of these audiences whichever level it chose. SmartLoad makes
the level of autonomy a setting: the safe-mode kill switch and the shadow-mode
default let a cautious team adopt the same product as an aggressive one, each at the
level of delegation they are prepared to defend to their own management. The cultural
impact, then, is not that SmartLoad imposes a new way of working, but that it
\emph{accommodates} the range of ways teams already feel about handing authority to
software --- which is what lets it be adopted across organisations with very
different cultures of caution.

\section{Identified Major Impact Across the Four Contexts}
\label{disc:synthesis}

Drawing the preceding sections together, \cref{tab:disc:contexts} states the single
most significant impact SmartLoad has --- or is positioned to have --- in each of the
four contexts the project is assessed against. The table is deliberately a synthesis
of major impacts rather than an exhaustive list: each row names the one effect that
most defines the project's significance in that context, together with the evidence
or boundary that keeps the claim grounded.

\begin{table}[ht]
  \centering
  \caption{Major identified impact of SmartLoad across the four assessment contexts.}
  \label{tab:disc:contexts}
  \small
  \begin{tabularx}{\linewidth}{@{}p{2.4cm}Yp{3.7cm}@{}}
    \toprule
    \textbf{Context} & \textbf{Major identified impact} & \textbf{Evidence / grounded boundary} \\
    \midrule
    Global / economic
      & Brings proactive, right-sizing adaptive traffic management to the
        sub-orchestration team, removing the over-provisioning premium and the
        per-node licensing cost without requiring a Kubernetes control plane.
      & Demonstrated bidirectional autoscaling; cost saving not yet quantified
        in money (\TBD{quantify monetary saving from a representative workload}). \\[2pt]
    Environmental
      & Reduces idle compute --- and therefore energy and carbon --- by scaling the
        pool down during demand troughs as a default behaviour, not only up at peaks.
      & Bidirectional scale-down shown in the committed run; per-deployment saving
        is modest and subject to the rebound effect. \\[2pt]
    Societal
      & Makes the online services people depend on more reliable at the highest-stake
        moments, while bounding the risk of opaque automation through a deterministic
        fallback and operator override on every decision.
      & Safety architecture verified; reliability benefit argued from mechanism, with
        a production deployment as the next step. \\[2pt]
    Cultural
      & Shifts operator work from reactive firefighting toward supervised automation,
        and lets teams with differing appetites for autonomy adopt the same tool via a
        kill switch and shadow-mode default.
      & Transparency surfaces (live feed, audit log) and the safe-mode flag are
        shipped and operable. \\
    \bottomrule
  \end{tabularx}
\end{table}

The common thread across all four rows is the project's defining commitment: a
system that places learned automation in the request path while guaranteeing that
the loss of that intelligence reduces it to classical infrastructure and never to an
outage. That commitment is what makes the economic case adoptable, the environmental
saving safe to enable by default, the societal benefit defensible against the failure
modes of opaque automation, and the cultural accommodation possible across teams that
feel very differently about trusting software with authority. The performance picture
of \cref{disc:performance} --- a sound mechanism with learned components at differing
maturity --- and the impact picture of this section are the same story told twice:
SmartLoad's significance lies not in any single model's score but in an architecture
disciplined enough to ship learned intelligence and classical safety in the same
system, and to be clear about which is carrying the load today.
